\chapter{怎样认出一只猫}
% 完整版：chapters/12_recognition.tex

\pencilfig{12}{\pencilart{12}}{大猫和小猫在眼睛的传感器上是截然不同的图案，答案却是同一个。}

图片里的猫可能在左边或右边，近处或远处，阳光下或阴影里。每一次，落在眼睛传感器上的都是完全不同的图案，而你在几分之一秒内就认出了它。这正是识别的难处：同一样东西，要在成千上万种不同的样子里认出来。

\section*{流水线}

大脑用一条大约十级的流水线来解决这个问题，信号走完全程只需十分之一点五秒。每一级做两件事。先是“与”：当需要的细节按需要的方式排列时，神经元就触发，比如两条线段相交成一个角。然后是“对所有位置的或”：只要附近任何地方有这个角，下一个神经元就触发，不管具体在哪儿。第一步让识别变得精细，第二步让它容忍平移。一级接一级，细节组合成部件，部件组合成物体，平移和缩放都被原谅了。

\pencilfig{112}{%
\foreach \i/\y/\cx/\cy in {1/1.2/0.15/0.3,2/0.3/0.3/0.1,3/-0.6/0.05/0.05}{
  \draw[penlite] (0,\y-0.3) rectangle (0.6,\y+0.3);
  \draw[pcGray, line width=0.8pt] (\cx,\y-0.3+\cy+0.35) -- (\cx,\y-0.3+\cy) -- (\cx+0.3,\y-0.3+\cy);
  \draw[pen=pcBlue, wash=pcBlue] (1.9,\y) circle (0.2);
  \draw[arr] (0.65,\y) -- (1.65,\y);
  \draw[arr=pcBlue] (2.12,\y) -- (3.62,{0.3+0.12*(2-\i)});}
\draw[pen=pcBlue, wash=pcBlue] (3.9,0.3) circle (0.25);
\draw[arr] (4.2,0.3) -- (5.75,0.3); \node[lbl, anchor=south] at (4.97,0.35) {更多级};
% a small cat
\draw[penlite] (6.3,0.3) circle (0.32);
\draw[penlite] (6.03,0.5) -- (6.07,0.82) -- (6.23,0.6); \draw[penlite] (6.57,0.5) -- (6.53,0.82) -- (6.37,0.6);
\node[lbl, anchor=north] at (0.3,-1.0) {图片}; \node[lbl, anchor=north, align=center] at (1.9,-1.0) {与：角\\在此处};
\node[lbl, anchor=north, align=center] at (3.9,-1.0) {或：角\\在任意处}; \node[lbl, anchor=north] at (6.3,-1.0) {猫};
}{流水线的一级：“与”在一个位置上拼出一个角，“或”不再计较它具体在哪儿。}

手机里识别图片的卷积神经网络，复制的正是这个方案。它们也投桃报李：训练来识别物体的网络，能相当好地预测视觉高层各级的神经元如何响应。

\section*{向视频学习}

可是，谁教过大脑？没有人给你看过一百万张标着“猫”的图片。一个可信的答案是：视频。世界是连续流动的，在几分之一秒内，物体本身几乎不变，变的是它的位置和光照。如果此刻看到的东西和片刻之前看到的东西之间的连接会增强，网络就会自己把同一物体的不同样子联系起来：变化慢的是物体，变化快的是环境。有一个实验在眼球运动的瞬间悄悄调换物体——神经元便开始把不同的物体当成同一个。作为一条普遍规则，这目前还是假说。

\section*{错觉是机制留下的指纹}

错觉不是故障，而是机器构造留下的痕迹。如果眼睛长时间盯着瀑布，再看静止的石头，石头会“往上爬”：“向下”和“向上”这一对通道中的一个疲劳了，两者失去平衡。那条著名的裙子，有人看成白金色，有人看成蓝黑色，它按照大脑默默假定的光照把人分成了两派。画出来的“旋转的蛇”看上去在动，是因为明亮区域的处理比暗淡区域稍快一点，时间差被读成了运动。而那个一会儿凸出、一会儿凹进的立方体，是两群神经元在相互竞争；按照一种模型，结果的切换主要由噪声决定。

有意思的是，一个训练来预测视频下一帧的神经网络，也“看见”蛇在旋转。可见，有些错觉是会预测所必须付出的代价。

\section*{大脑对神经网络}

两者并不完全相像。大脑只需少量例子，几乎不需要标注就能学习，而且做这一切只用二十瓦。神经网络可以被人眼察觉不到的图片篡改所欺骗。不过，这类篡改如果一闪而过，也会让人稍稍看走眼。这两台机器之间真正的界线在哪里，还在研究之中。

\fullversion{12}
